\documentclass[a4paper,11pt]{book}

\usepackage[T1]{fontenc}
\usepackage[utf8]{inputenc}
\usepackage{graphicx}
\usepackage{xcolor}
\usepackage[colorlinks=true, linkcolor=blue]{hyperref}

\usepackage{mathptmx}
\usepackage{amsmath,amssymb,amsthm}
\usepackage{bbm}

\usepackage{geometry}
\geometry{left=15mm, right=15mm, top=20mm, bottom=20mm}

% Theorem style
\makeatletter
\def\th@plain{%
  \thm@notefont{}%
  \itshape
}
\def\th@definition{%
  \thm@notefont{}%
  \normalfont
}
\makeatother

% Numbering
\renewcommand{\thechapter}{\Roman{chapter}}
\renewcommand*\thesection{\arabic{section}}

% Theorem environments
\newtheorem{theorem}{Theorem}[section]
\renewcommand{\thetheorem}{\arabic{section}.\arabic{theorem}}
\newtheorem{corollary}{Corollary}[theorem]
\newtheorem{lemma}[theorem]{Lemma}
\newtheorem{proposition}[theorem]{Proposition}

\theoremstyle{definition}
\newtheorem{definition}{Definition}[section]
\renewcommand{\thedefinition}{\arabic{section}.\arabic{definition}}
\newtheorem*{example}{Example}

\theoremstyle{remark}
\newtheorem*{remark}{Remark}
\newtheorem{review}{Review}

% Math operators
\DeclareMathOperator{\var}{Var}

\linespread{1.3}

\begin{document}
\title{Stochastic I}

\chapter{Probability Theory}
\section{Probability Distribution}
\begin{theorem}[Transformation Formula]
Given a random variable $X \in \mathbb{R}^n$ with joint density $f$, $g: \mathbb{R}^n \rightarrow \mathbb{R}^n$ a continuously differentiable and injective function, with non-vanishing Jacobian, then the density function of $Y = g(X)$ is:
\begin{align*}
	f^Y(y) = f^X(g^{-1}(y)) |\det J_{g^{-1}}(y)|
\end{align*}
\end{theorem}

\begin{example}
\end{example}

\section{Integration with Respect to probability measure}
\begin{theorem}[Markov Inequality]
Given a random variable $X$ and \textbf{non-negative} monotone increasing function $f:[0, \infty) \rightarrow [0, \infty)$, then we have:
\begin{align*}
	\forall \epsilon > 0 \textbf{ with } f(\epsilon) > 0: \mathbb{P}(|X| \geq \epsilon) \leq \frac{\mathbb{E}f(|X|)}{f(\epsilon)}
\end{align*}
\end{theorem}
\begin{proof}
	\begin{align*}
		\mathbb{E}[f(X)] &\geq \mathbb{E}[f(X) \mathbbm{1}_{f(X) \geq f(\epsilon)}]\\
		&\geq \mathbb{E}[f(\epsilon)\mathbbm{1}_{f(X) \geq f(\epsilon)}]\\
		&= f(\epsilon) \mathbb{P}(f(X) \geq f(\epsilon))\\
	(1)	&\leq f(\epsilon)\mathbb{P}(|X| \geq \epsilon)
	\end{align*}
	where in (1) we use that $f$ is an increasing function.
\end{proof}

\begin{proposition}[Fubini-Tonelli]
	Given $X$ a non-negative random variable, then:
	\begin{align*}
		\mathbb{E} X = \int_{\mathbb{R}} \mathbb{P}(X \geq t) \, dt
	\end{align*}
\end{proposition}
\begin{proof}
Observe:
	\begin{align*}
		\int_0^\infty \mathbb{P}(X(\omega) \geq x) \, dx &= \int_{\mathbb{R}^+} \int_\Omega \mathbbm{1}_{X(\omega) \geq x} \, d\omega \, dx\\
		&= \int_\Omega \int_0^{X(\omega)} dx \, d\omega \\
		&= \int_\Omega X(\omega) \, d\omega \\
	 	&= \mathbb{E} X
	\end{align*}
\end{proof}


\section{Terminal $\sigma$-algebra and independence}
This chapter we refer mainly to \textbf{Klenke}.
\subsection{Borel-Cantelli}
\begin{definition}[Limes superior and limes inferior]
	Given $(A_n)_{n \geq 1}$ a sequence of subsets in $X$. We define:
	\begin{align*}
		\limsup_n A_n &:= \{ x \in X: x \in A_n \textbf{ for infinitely many } A_n \}\\
		&\text{(sequence is decreasing)}\\
		\liminf_n A_n &:= \{ x \in X: \exists n_0(x) \in \mathbb{N}\,\forall n \geq n_0(x): x \in A_n ( \textbf{ x is in all but finitely many $A_n$})\}\\
		 &\text{(sequence is increasing)}
	\end{align*}
	equivalently we see:
	\begin{align*}
		\limsup_n A_n &= \bigcap_{n = 1}^\infty \bigcup_{k = n}^\infty A_n\\
		\liminf_n A_n &= \bigcup_{n = 1}^\infty \bigcap_{k = n}^\infty A_n
	\end{align*}
	We define the convergence of sets as:
	\begin{align*}
		\limsup_n A_n = \liminf_n A_n := \lim_n A_n
	\end{align*}
\end{definition}
\begin{theorem}[The first Borel-Cantelli Lemma]
	If $\sum_{n = 1}^\infty \mathbb{P}(A_n) < \infty$:
	\begin{equation*}
		\mathbb{P}(A_n \text{ i.o.}) = 0
	\end{equation*}
\end{theorem}
\begin{proof}
\begin{align*}
	\sum \mathbb{P}(A_n) < \infty &\Rightarrow \lim_n \mathbb{P}(A_n) = 0\\
	&\Rightarrow \mathbb{P}(\lim_n A_n) = 0\\
	&\Rightarrow \mathbb{P}(\limsup A_n) = 0\\
	&\Rightarrow \mathbb{P}(A_n \text{ i.o.}) = 0
\end{align*}
\end{proof}

\begin{theorem}[The second Borel-Cantelli Lemma]
	$A_n$ independent and $\sum P(A_n) = \infty$, then
	\begin{equation*}
		\mathbb{P}(A_n \text{ i.o.}) = 1
	\end{equation*}
\end{theorem}

\begin{theorem}
	$X_1, X_2, \dots$ i.i.d.\ with $\mathbb{E} X_i = \infty$, then
	\begin{equation*}
		\mathbb{P}(|X_n| \geq n \text{ i.o.}) = 1
	\end{equation*}
	So if $S_n = X_1 + \dots + X_n$ then:
	\begin{equation*}
		\mathbb{P}\!\left(\lim \frac{S_n}{n} \in (-\infty, \infty) \text{ i.o.}\right) = 0
	\end{equation*}
\end{theorem}


\subsection{Kolmogorov 0-1 Law}
\begin{definition}[Tail $\sigma$-algebra]
Given a set of $\sigma$-fields $(\mathfrak{A}_i)_{i \in I}$. The \textbf{tail $\sigma$-field} is defined as:
	\begin{equation*}
		\mathfrak{T}\!\left((\mathfrak{A}_i)_{i \in I}\right) := \bigcap_{\substack{J \subset I \\ |J| < \infty}} \sigma\!\left(\bigcup_{j \in I \setminus J} \mathcal{A}_j\right)
	\end{equation*}
\end{definition}
\begin{remark}
Related to the limes superior of the sets.
\end{remark}


\begin{theorem}[Kolmogorov 0-1 Law]
	$I$ countable index set, $(\mathcal{A}_i)_{i \in I}$ independent $\sigma$-fields, then
	\begin{equation*}
		\forall A \in \mathfrak{T}\!\left((\mathfrak{A}_i)_{i \in I}\right): \mathbb{P}(A) \in \{0, 1\}
	\end{equation*}
\end{theorem}
\begin{proof}

\end{proof}

\begin{theorem}[Hewitt-Savage 0-1 Law]
	Define $\mathcal{E}$ the $\sigma$-field of all events stable under finite permutation (prove it's indeed a $\sigma$-field). Given $X_1, \dots, X_n, \dots$ i.i.d.\ random variables, if $A \in \mathcal{E}$, then $\mathbb{P}(A) \in \{0, 1\}$.
\end{theorem}


\subsection{Independence, Distribution and Expectation}
\begin{review}[Central limit theory]
The regularised mean of a sequence of random variables converges in distribution to the standard normal distribution (under certain conditions).
\begin{equation*}
    Z := \frac{\sqrt{n}}{\sigma}(\overline{X} - \mu) \overset{d}{\longrightarrow} \mathcal{N}(0, 1)
\end{equation*}
\end{review}

\begin{review}[Slutsky lemma]
If a sequence of random variables converges in probability to a constant while another converges in distribution to a random variable, then the linearity holds (the linear combination converges in distribution).
\end{review}

\begin{review}[Portmanteau theorem]
\end{review}

\section{Convergence Concept of Random Variables}
\begin{definition}[Convergence in probability]
\end{definition}

\begin{definition}[Convergence in distribution]
\end{definition}

\begin{definition}[Convergence in $L^p$]
\end{definition}



\section{Law of Large Numbers}
\begin{theorem}[Strong law of large numbers]
	Let $(X_n)_{n \ge 1}$ be a sequence of i.i.d.\ random variables such that $X_j \in L^2(\Omega)$ with
	\begin{align*}
		\mathbb{E}X_j = \mu, \quad \var X_j = \sigma^2 \qquad \forall j = 1, 2, \dots
	\end{align*}
	Define the running sum $S_n := \sum_{j = 1}^n X_j$, then
	\begin{align*}
		\frac{S_n}{n} \overset{a.s.}{\longrightarrow} \mu \quad \text{and} \quad \frac{S_n}{n} \overset{L^2}{\longrightarrow} \mu.
	\end{align*}
\end{theorem}


\begin{theorem}[$L^2$ weak law]
	$X_1, X_2, \dots$ uncorrelated with $\mathbb{E}X_i = \mu$ and $X_i \in L^2$, then:
	\begin{equation*}
		\frac{S_n}{n} \overset{L^2}{\longrightarrow}\mu, \qquad \frac{S_n}{n} \overset{\mathbb{P}}{\longrightarrow} \mu.
	\end{equation*}
\end{theorem}

\begin{lemma}
	If $Y\geq 0$ and $p >0$, then:
	\begin{equation*}
		\mathbb{E} Y^p = \int_{0}^\infty p y^{p-1} \mathbb{P}(Y > y) \, dy
	\end{equation*}
\end{lemma}



\section{Central Limit Theory}
\subsection{Characteristic functions}
\begin{definition}[Characteristic function]
	Given a random variable $X \in L^1(\mathbb{P})$. The characteristic function is the Fourier transform of the distribution of $X$:
	\begin{align*}
		\varphi(u): \mathbb{R}^n \rightarrow \mathbb{C}, \quad u \mapsto \mathbb{E}[e^{iu \cdot X}]
	\end{align*}
\end{definition}

\begin{example}[The normal distribution]
\end{example}



\begin{theorem}[The inversion formula]
	Let $\varphi(t)$ be the characteristic function of the probability measure $\mu$ of the random variable $X$. If $a<b$, then:
	\begin{equation*}
		\lim_{T \rightarrow \infty} (2\pi)^{-1} \int_{-T}^T \frac{e^{-ita} - e^{-itb}}{it} \varphi(t) \, dt = \mu(a, b) + \frac{1}{2} \mu(\{a, b\})
	\end{equation*}
\end{theorem}

\begin{theorem}
	If $\int |\varphi(t)| \, dt < \infty$, then $\mu$ has the bounded continuous density:
	\begin{equation*}
		f(y) = \frac{1}{2\pi} \int e^{-ity} \varphi(t) \, dt
	\end{equation*}
\end{theorem}

\begin{theorem}[Continuity theorem of L\'evy]
	Let $\mu_n$, $1\leq n \leq \infty$ be probability measures with characteristic functions $\varphi_n$, then:
	\begin{enumerate}
		\item $\mu_n \rightarrow \mu_\infty \quad \Rightarrow \quad \varphi_n(t) \rightarrow \varphi_\infty(t) \quad \forall t$
		\item $\varphi_n(t) \rightarrow \varphi(t)$ for some $\varphi$ and $\varphi$ is continuous at $0$, then the associated sequence of distributions $\mu_n$ is tight and converges weakly to the measure $\mu$ with characteristic function $\varphi$.
	\end{enumerate}
\end{theorem}

\begin{definition}
Let $X$ be a $\mathbb{R}^n$ random variable, $k \in \mathbb{R}^n$, then we define the characteristic function of $X$ as:
	\begin{equation*}
		\varphi^X(k) := \int_{\mathbb{R}^n} e^{i k \cdot x} \, \mathbb{P}^X(dx)
	\end{equation*}
\end{definition}

\begin{theorem}[Uniqueness theorem]

\end{theorem}

\begin{theorem}[Central limit theorem I]
	Given $X_1, \dots, X_n, \dots$ a sequence of i.i.d.\ samples such that $X_i \in L^2(\Omega)$ with $\mathbb{E}X_i= \mu$, $\var X_i = \sigma^2$. Define the running sum $S_n:= \sum_{i = 1}^n X_i$ and $Y_n:= \frac{S_n - n \mu}{\sqrt{n} \sigma}$. Then
	\begin{align*}
		Y_n \overset{d}{\longrightarrow} \mathcal{N}(0, 1).
	\end{align*}
\end{theorem}

\begin{theorem}[Central limit theorem II]
Given $X_1, \dots, X_n, \dots$ a sequence of independent (but not necessarily identically distributed) samples such that $X_i \in L^2(\Omega)$ with $\mathbb{E}X_i= 0$, $\var X_i = \sigma^2_j$. Define the running sum $S_n:= \sum_{i = 1}^n X_i$. If
\begin{align*}
	&\sup_j \mathbb{E} |X_j|^{2 + \epsilon}< \infty \text{ for some } \epsilon >0, \\
	&\sum_{j = 1}^\infty \sigma_j^2 = \infty,
\end{align*}
then
\begin{align*}
	\left(\sqrt{\sum_{j = 1}^n \sigma_j^2}\right)^{-1} S_n \overset{d}{\longrightarrow} \mathcal{N}(0, 1).
\end{align*}
\end{theorem}

\section{The Gaussian Random Variable}



\section{Weak Convergence}
\begin{definition}[Separating Family]
	Given $F \subset M(E)$ a subset of Radon measures (inner regular Borel measures). A family of measurable functions $C$ is a separating family for $F$ if $\forall \mu, \nu \in F$:
	\begin{equation*}
		\int f \, d\mu = \int f \, d\nu \quad \forall f \in C \cap \mathcal{L}^1(\mu) \cap \mathcal{L}^1(\nu) \quad \Rightarrow \quad \mu = \nu
	\end{equation*}
\end{definition}

\begin{definition}[Weak convergence of measures]

\end{definition}


\begin{definition}
	A sequence of random variables converges weakly or converges in distribution if its distribution function converges for all continuity points of $F$.
\end{definition}

\begin{theorem}
	If $F_n \rightarrow F_\infty$
\end{theorem}

\begin{theorem}
	Every probability measure on a Polish space is tight.
\end{theorem}

\begin{theorem}
	Every probability measure on a Polish space is regular.
\end{theorem}

\begin{theorem}[Continuous Mapping theorem]
	If $g: S \rightarrow T$ is continuous, $S$, $T$ metric spaces, then:
	\begin{equation*}
		X_n \overset{d}{\rightarrow} X \Rightarrow g(X_n) \overset{d}{\rightarrow} g(X)
	\end{equation*}
\end{theorem}
\begin{proof}
	Notice that $\forall f \in C_b(T, \mathbb{R})$, $f\circ g$ is also bounded continuous.
\end{proof}
\begin{remark}
If $X_n \overset{a.s.}{\longrightarrow} X$, then $g(X_n) \overset{a.s.}{\longrightarrow} g(X)$.
\end{remark}

\begin{theorem}
	Convergence in distribution implies convergence of the cdf at every \textbf{continuity} point.
\end{theorem}

\end{document}
